\section{Constructing charts, residue functionals and Hermitian carriers}
\label{sec:global-forms}

\status{FC1--FC18 prove the arithmetic and linear-algebra constructions
specified below, including actual finite charts from supplied adelic
representatives, replacement of their residue functionals without changing
their level groups, exact residue projections, the Fock metric and period
identities, the common auxiliary-prime data, and the Hermitian carrier
conversion. FC19 joins these with CV9 and NB7--NB8 to construct the
actual step-zero \texttt{Forms0}, its matching level carrier, and all
the global input records needed by VE9. This construction retains the
initial decorated-point realizability input. The complete simultaneous
\texttt{FormsTiedSym} and \texttt{ColsTiedSym} constructors still require
the terminal incidence construction and the specified symmetric
auxiliary-prime level descriptions. The original \texttt{TC.Ok},
\texttt{Pins}, \texttt{Laws} and \texttt{Later} data remain hypotheses.}

\paragraph{The data used here.}
Write $E/F$ for the supplied CM extension, $\mathcal O=\mathcal O_E$,
and $x\mapsto\bar x$ for its involution. Fix a word $u$, a step $i\le m$,
and $n=N_i>0$. The exact coordinate identifications of SF2 put the
Hermitian form in the form $h(x,y)=\bar x^{\mathsf T}y$ and preserve the
original class lattice and frame matrices. Put $p=T_u$ and
$\pi=\pi_u$, so $p$ is a rational prime and
$\pi\bar\pi=p$. The prime data include
$(\pi,\bar\pi)=\mathcal O$, the selected set
$S=\{P:P\mid\pi\}$, $|S|=[F:\mathbb Q]$, and
$|\mathcal O/P|=p$ for $P\in S$.

The source's \texttt{Laws} are substantive supplied data. They include
finite orbit and stabilizer statements for the actual column actions,
the level and period action laws, and the polynomial lattices and their
pairings. They do not mean merely that some unnamed compatible objects
exist. Whenever the proof below uses adelic realizability or a supported
move not yet reconstructed, that use is stated separately.

\paragraph{FC1: the level-prime index and its real residue field.}
First, $(\pi)=\prod_{P\in S}P$. Indeed, conjugation preserves absolute
ideal norm, and $[E:\mathbb Q]=2[F:\mathbb Q]=2f$, so
$\mathrm N((\pi))^2=p^{2f}$ and $\mathrm N((\pi))=p^f$.
The $f$ distinct primes in $S$ already contribute $p^f$; none can have
a larger exponent. Since distinct maximal ideals are comaximal,
\[
 r\in\bigcap_{P\in S}P\quad\Longleftrightarrow\quad\pi\mid r.
 \tag{FC.1}
\]
This is the \texttt{BasePi.dvd} assertion, with its exact squarefree
content, rather than merely a statement about the support of $(\pi)$.

Restriction gives a bijection from $S$ to the primes $v$ of $F$ above
$p$. For surjectivity, choose a prime $Q$ of $E$ over $v$. Since
$p=\pi\bar\pi\in Q$, either $\pi\in Q$, or
$\pi\in\bar Q$; conjugate primes have the same restriction. For
injectivity, the primes over a fixed $v$ in this quadratic Galois
extension are one conjugate orbit. Two selected primes cannot be
conjugate, since that would put both $\pi$ and $\bar\pi$ in one
prime, contrary to their coprimality. Thus the map is the literal
\texttt{underT}; its inverse does not select a different orientation.

There is also an explicit residue-field descent. Let
$\rho:\mathcal O\twoheadrightarrow k$ be the chosen quotient at $P$,
where $\rho(\pi)=0$ and $\rho(\bar\pi)\ne0$, and let
$v=P\cap\mathcal O_F$. For $y\in\mathcal O$, the elements
\[
 t=\pi+\bar\pi,\qquad z=y\bar\pi+\bar y\pi
 \tag{FC.2}
\]
are integral and fixed by conjugation, hence belong to $\mathcal O_F$.
The residue of $t$ is nonzero. Choose $s\in\mathcal O_F$ with
$st\equiv1\pmod v$. Then $\rho(sz)=\rho(y)$. Consequently the
restriction of $\rho$ to $\mathcal O_F$ is onto, has kernel $v$, and
induces the exact isomorphism
\[
 k\simeq\mathcal O_F/v,
 \qquad \rho(a)\longmapsto a\bmod v\quad(a\in\mathcal O_F).
 \tag{FC.3}
\]
No choice of an abstract field of the same cardinality is substituted
for this map. Together with the supplied embedding/prime correspondence,
these are the prime-index and residue fields of \texttt{atTTied} and
\texttt{CarH.fv}. The coefficient-place labels \texttt{tauOf} and
\texttt{seenP} still use that supplied arithmetic correspondence.

\paragraph{FC2: coupling data are determined by their reductions, and
every prescribed reduction can be lifted.}
Here the reference twist-prime set is empty, as it is for $\mu_0$ in
the source. For a decorated lattice $D$, write $D_P$ for its localized
lattice and $b_P:D_P\to k_P^n$ for its supplied reduction frame.
At an auxiliary prime $P$, the rational integer $q_P$ lies in
$P\setminus P^2$; thus it is a uniformizer in $\mathcal O_P$ and
\[
 b_P(x)=0\quad\Longleftrightarrow\quad x\in q_P D_P
 \quad(x\in D_P).                                      \tag{FC.4}
\]
This follows by writing $x$ in the local frame: zero reduction means
that every coordinate is divisible by the uniformizer.

At the coupling step $i=m-1$, the additional datum is a coset
$y+M$, where $q_Py\in D_P$ and $M$ imposes $x\in D_P$ at each
auxiliary prime. Its reduction at $P$ is the singleton
$\{b_P(q_Py)\}$. If two such cosets have the same reductions,
(FC.4) gives $y-y'\in D_P$ at every relevant $P$, hence
$y-y'\in M$ and the cosets agree. At all other steps the additional
set is empty. Thus a column point is determined by its decorated
lattice and the family of coupling reductions. This proves the actual
\texttt{eq\_of\_deco\_bar} assertion in the reference-tuple case.

For existence, prescribe vectors $v_P\in k_P^n$ for an arbitrary
decorated lattice $D$. Its global reduction frames supply
$x_P\in D$ with $b_P(x_P)=v_P$. Choose $e_P\in\mathcal O$ with
$e_P\equiv1\pmod P$ and such that $e_P/q_P$ is integral at every
other auxiliary prime. Such elements follow from the Chinese remainder
theorem, using sufficiently large powers at the other primes. Then
\[
 y=\sum_P\frac{e_P}{q_P}x_P                            \tag{FC.5}
\]
satisfies $q_Qy\in D_Q$ and $b_Q(q_Qy)=v_Q$ for every $Q$:
the $Q$-summand gives $v_Q$, and each other summand is a multiple of
$q_Q$ in $D_Q$. The coset $y+M$ is therefore the required point.
This proves \texttt{exists\_ptAt\_bar\_eq}; no adelic approximation
theorem is being used in this finite residue construction.

\paragraph{FC3: the finite chart fibers and their exact representatives.}
Let $\mathscr C_i$ be the supplied finite column-class set. A class
$c$ has an index $b(c)$ and a preferred point $r(c)$; let
$C(c)$ denote its coupling-reduction family. Define
\[
 \operatorname{Ch}_i=\operatorname{im}
       \bigl(c\mapsto(b(c),C(c))\bigr),\qquad
 \operatorname{chOf}(c)=(b(c),C(c)).                     \tag{FC.6}
\]
This is a finite image, not a new finite-class assumption. Choose a
representative $c_\beta$ of each image point. If $H_b$ is the original
product of finite reduction groups, put
$H_{b,C}=H_b\cap\operatorname{Stab}(C)$, component by component.

The fiber over $\beta=(b,C)$ is exactly the decorated-lattice orbit
set for $\mathrm U(E)\times H_{b,C}$. To prove injectivity, an element
of this group relating the decorated lattices also preserves $C$;
FC2 then says that it relates the whole points. They represent the same
original class. For surjectivity, lift any decorated lattice with the
prescribed $C$ by (FC.5), and take its preferred representative. Preferred
representatives normalize the coupling coset by choosing one value from
its orbit. Starting with the same $C$ therefore gives the same normalized
value: the two sets of attainable values are identical. The moving
element consequently lies in $H_{b,C}$. This proves both directions,
including the normalization step which would be missing from an arbitrary
partition into fibers.

For the passage to adelic charts we retain one precise geometric input:
each actual decorated lattice has a realizing unitary adelic point.
AD1 constructs the diagonal base point $k_u$ with the lattice and
last-row residues specified in \texttt{kStd\_spec}. Choose the actual source representative
$y_c=y_{\rm Std}(u,i,c)$ so that $y_ck_u$ realizes $r(c)$.
Pulling a vector back by $y_c^{-1}$ gives the chart base lattice
$L_\pi=\mathcal O^{n-1}\oplus\pi\mathcal O$. Equality of the
decorated points says exactly that a quotient matrix preserves this
lattice and has reduction in $H_{b,C}$: the lattice equality checks
integrality of the matrix and its inverse, and equality of the finite
decorations checks the reduction condition. Rational changes of frame
act on the left. Thus the fiber bijection just proved is also the
adelic double-quotient bijection required of a \texttt{LevelChart}.
The supplied injective stabilizer map has exactly the rational elements
in $y_cUy_c^{-1}$ as its image, by the same equality-of-points criterion.

The chart character is the restriction of the original finite character,
composed with reduction. More explicitly, right multiplication by an
adelic $x\in U$ acts on the decoration by the inverse of its reduction.
The definition of the level character contains this same inverse, so
the two inversions cancel. Its value is precisely the original column
character of the moving element. This proves the character clause of
\texttt{ReadsK}, in addition to its lattice, level and stabilizer clauses.
The matrices remain $y_{\rm Std}$, not representatives chosen anew
inside each chart. At $i=0,m$ the column index is unique and the coupling
set is empty; the nonempty image (FC.6) is therefore a singleton.
This gives the exact endpoint \texttt{Unique} assertions.

\paragraph{FC4: normalize the residue functional without changing the
level.}
The following local calculation supplies the substantial part of
\texttt{charts\_funKilled}. Let $A$ be the integral algebra
$\mathcal O_E\otimes_{\mathcal O_F}\mathcal O_{F_v}$ inside
$E\otimes_FF_v$, with its involution; write $q=\bar\pi$.
At a selected level prime, (FC.1) and FC1 identify both $A/\pi A$
and $A/qA$ with fields of order $p$. Put
\[
 L=A^{n-1}\oplus\pi A,\quad
 L^\#=A^{n-1}\oplus q^{-1}A,\quad
 L^\#/L\simeq A/pA,\quad [x]\longmapsto[qx_n].         \tag{FC.7}
\]
The middle identity is SF3, after base change. Suppose $g$ is unitary
and preserves $L$. It also preserves $L^\#$. Applying $g$ to
$\pi e_j\in L$ shows that its last row is integral; applying it to
$e_j\in L^\#$ shows that all other rows are integral. Moreover
\[
 g_{nj}\in\pi A\ (j<n),\qquad
 g_{jn}\in qA\ (j<n).                                 \tag{FC.8}
\]
On the quotient (FC.7), $g$ therefore acts by multiplication by
$d=g_{nn}$. The $(n,n)$ entry of $g^*g=1$ gives
\[
 d\bar d\equiv1\pmod{pA}.                             \tag{FC.9}
\]

Let $f:L^\#\to M$ be any nonzero $A$-linear functional vanishing on
$L$, where the additive group of $M$ has prime order $p$. Multiplication
by $\pi$ on $M$ is either zero or onto. In the second case
$qM=0$, since $q\pi=p$ and $pM=0$. Thus $f$ factors nontrivially
through either $A/\pi A$ or $A/qA$. The resulting map to $M$ is
an isomorphism: both additive groups have prime order. It follows that
$f(gx)=f(x)$ for every $x$ exactly when $d\equiv1$ modulo the
chosen one of these ideals. By (FC.9) and conjugation, the two conditions
are equivalent. In particular they are also equivalent to $g$ acting
trivially on all of $L^\#/L$.

Replace $f$ by the canonical quotient
\[
 f_0:L^\#\longrightarrow A/\pi A,\qquad
 f_0(x)=qx_n\bmod\pi A.                               \tag{FC.10}
\]
It is nonzero, kills $L$, has a value group of order $p$, and
$f_0(\pi x)=0$. The intersection of its stabilizer with the unitary
stabilizer of $L$ is exactly the old intersection. Therefore one can
replace only the chart's value module and functional: its base lattice,
every local level group, global level, character, adelic representatives
and stabilizer maps stay literally unchanged. Its \texttt{ReadsK}
proof stays valid. Doing this independently at each level prime and
each chart yields the actual \texttt{FunKilled} property. For $n=1$
the off-diagonal statements are empty and the same proof applies.

\paragraph{FC5: the quotient projections are constructed, not assumed.}
Let $\mathfrak p\subset\mathcal O$ be maximal and let two lattices
$A_0,M_0\subset E^n$ satisfy $(A_0)_{\mathfrak p}=(M_0)_{\mathfrak p}$.
For $x\in A_0$, choose $s\notin\mathfrak p$ with $sx\in M_0$ and set
\[
 \operatorname{pr}(x)=\bar s^{-1}[sx]
                 \in M_0/\mathfrak pM_0.              \tag{FC.11}
\]
Using a common denominator proves independence of $s$, additivity and
$\operatorname{pr}(ax)=\bar a\operatorname{pr}(x)$.
If $x\in M_0$, this is its usual residue. Given $[y]\in M_0/\mathfrak pM_0$,
choose $s\notin\mathfrak p$ with $sy\in A_0$ and $t$ with
$ts\equiv1\pmod{\mathfrak p}$; then $tsy\in A_0$ maps to $[y]$.
Thus $\operatorname{pr}$ is onto. Its kernel is exactly
$\mathfrak pA_0$: if its value is zero, clearing a second denominator
gives $rx\in\mathfrak pA_0$ for some $r\notin\mathfrak p$;
choose $a$ with $ar\equiv1\pmod{\mathfrak p}$ and write
$x=arx+(1-ar)x$.

For the actual decoration, its defining local laws give
$(L^\#)_{\mathfrak p}=N_{\mathfrak p}$ and
$(L+\mathfrak pL^\#)_{\mathfrak p}=(\operatorname{dn}N)_{\mathfrak p}$.
Apply (FC.11) to these two pairs. They are the source's
\texttt{pr} and \texttt{prk}, with the required surjectivity and kernels.
The decoration functional satisfies
\[
 f_{\mathfrak p}(z)=0\quad\Longleftrightarrow\quad
 z\in L+\mathfrak pL^\#\qquad(z\in L^\#).              \tag{FC.12}
\]
For the nontrivial implication, $f_{\mathfrak p}(z)=0$ puts $z$ in
$L_{\mathfrak p}$; clear a denominator $s$ and choose
$ts+r=1$, $r\in\mathfrak p$, to write
$z=t(sz)+rz\in L+\mathfrak pL^\#$. The converse follows by
linearity. Finally, if $f_{\mathfrak p}(z)=f_{\mathfrak p}(e)=1$,
then $z-e\in L_{\mathfrak p}$ and
$pz-pe\in\mathfrak p(\operatorname{dn}N)_{\mathfrak p}$.
Their projections agree. This is the special \texttt{te} compatibility,
and completes \texttt{exists\_level\_projections} with the literal
\texttt{IdL} identities on global elements.

\paragraph{FC6: the residue comparison has an explicit formula.}
Let $f:L^\#\to M$ now be a nonzero normalized functional, so
$f(\pi x)=0$. Its factorization in FC4 gives a unique additive
isomorphism $e:M\simeq k_P$ satisfying
\[
 e(f(x))=\rho_P(\bar\pi x_n).                          \tag{FC.13}
\]
Indeed, the map $x\mapsto\rho_P(\bar\pi x_n)$ is onto, and its
kernel is killed by $f$: subtract
$\pi\bar\pi^{-1}ae_n$ when $\bar\pi x_n=\pi a$; the
remainder lies in $L$. The induced nonzero map $k_P\to M$ is an
isomorphism between groups of order $p$.

This normalization is compatible with the actual decorated-class
functional, not merely with an abstract residue field. If $y_ck_u$
realizes the point, its inverse is $k_u^{-1}y_c^{-1}$. The defining
last-row residue of $k_u^{-1}$ is
$(0,\ldots,0,\bar\pi)$. Reduction of its product with
$y_c^{-1}z$ therefore gives
\[
 f_P(sz)=\bar s\,\rho_P\!\left(
        \bar\pi (y_{c,v}^{-1}z)_n\right)               \tag{FC.14}
\]
whenever $sz$ is in the original decorated lattice $N$.
All the products in this reduction are integral at $P$: the first
$n-1$ coordinates of $y_{c,v}^{-1}z$ are integral, and its last is
integral at $P$ because $\bar\pi$ is a unit there. This checks the
integrality needed when multiplying residue identities.

Choose $s\notin P$ with $sz\in N$, as permitted by FC5.
The defining rule (FC.11), (FC.13) and (FC.14) then agree after
multiplication by the nonzero scalar $\bar s$; cancel it.
This proves \texttt{FvTieG}. For a rational base-coordinate vector
$x$, it also proves the precise terminal formula
$e(f(x))=\operatorname{res}(r)$ when $r=\bar\pi x_n\in\mathcal O$.
Combined with $h(\pi e_n,x)=\bar\pi x_n$, this is the
\texttt{TilBaseFv} calculation once the terminal vector and its
incidence data have been supplied.

\paragraph{FC7: the metric comparison is a finite polynomial calculation.}
For polynomials in all frame entries write
$P=\sum_\alpha p_\alpha X^\alpha$ and define the normalized Fock form
\[
 N_\lambda(P,Q)=c_\lambda^{-1}
       \sum_\alpha\alpha!p_\alpha\overline{q_\alpha},
 \qquad c_\lambda>0,                                  \tag{FC.15}
\]
where $c_\lambda$ is the product of the Fock squared norms (the
self-pairings) of the source's highest polynomials. This convention is linear in $P$ and conjugate
linear in $Q$. It is positive definite. For completeness, the unnormalized
sum equals the Gaussian integral
$\pi^{-d}\int_{\mathbb C^d}P(z)\overline{Q(z)}e^{-|z|^2}\,dz$:
angular integration kills unequal exponents and radial integration gives
$\alpha!$. A unitary change of variables preserves this integral and
therefore the form. Blockwise left multiplication of the frame matrices
by unitary matrices is such a change of variables. Adding a common
integer to every entry of a highest weight changes neither its
normalized weight nor $c_\lambda$, hence does not change (FC.15).

In the actual standard coordinates, an orthonormal frame $o$ satisfies
$o_\tau^*o_\tau=I$, so $|\det o_\tau|=1$. Inserting the
determinant-twisted polynomial definitions gives
\[
 \operatorname{ipV}(f_P,f_Q)=N_\lambda(Q,P),
 \qquad f_P(z)=\operatorname{fnOf}(o^{-1},P)(z).
 \tag{FC.16}
\]
The order $Q,P$ is essential: $\operatorname{ipV}$ is conjugate linear
in its first argument. No reversal of the Hermitian convention is
being made.

The global Gram matrix is defined by the finite class sum with weights
$a_c=|\operatorname{Stab}(c)|^{-1}>0$. Expanding a vector in its
chosen basis and distributing this finite sum proves, at an embedding
$s$ of the coefficient field,
\[
 \operatorname{maj}_s(x,x)
   =\sum_c a_c\,N_\lambda(sP_c,sP_c)
   =\sum_c a_c\prod_\tau|\det o_\tau|^{2\delta_\tau}
       N_\lambda\bigl(P_c^s(oX),P_c^s(oX)\bigr).
 \tag{FC.17}
\]
Taking real parts is the source's \texttt{NormTieG}. The following
descent argument justifies the embedded Gram expansion even when the
coefficient field $L$ is larger than a CM field. Define
\[
 L^{\mathrm c}=\{a\in L:
    s(a^*)=\overline{s(a)}\text{ for every }s:L\hookrightarrow\mathbb C\}.
 \tag{FC.17a}
\]
It is a subfield, since the displayed equality is preserved by field
operations. It contains the subfield
$L_0=\mathbb Q(\tau(E),\chi(g):\tau,g)$ generated by the embedded
CM-field images and the character values in the stabilizer actions.
For the first generators this follows from the supplied
\texttt{FOk.emb\_star} and CM conjugation laws. For a root of unity
$\zeta$, its distinguished embedding has modulus one, and
\texttt{FOk.sig\_star} identifies star there with complex conjugation;
injectivity
of that embedding gives $\zeta^*=\zeta^{-1}$. Every other embedding
also maps it to a root of unity, hence
$s(\zeta^*)=s(\zeta)^{-1}=\overline{s(\zeta)}$.
These are the generators actually used by the polynomial action.

The integer Weyl-polynomial basis identifies the stabilizer equations
with linear equations over $L^{\mathrm c}$. Their invariant space has
an $L$-basis with coordinates in $L^{\mathrm c}$. One direct proof,
valid simultaneously for all equations, chooses an
$L^{\mathrm c}$-basis $a_j$ of $L$ and applies its coordinate
projections to every coefficient of an invariant vector. Each projection
commutes with every action matrix, so gives another invariant vector
with coefficients in $L^{\mathrm c}$. The original vector is
$\sum_j a_j$ times these projected vectors. The latter therefore span
over $L$, and a basis can be extracted from them.

Crucially, this is the basis used by the actual source: the definition
\texttt{CSD.basis} selects a basis with these coordinates whenever
\texttt{HasRatBasis} holds, as just proved. It is not an arbitrary
$L$-basis followed by an unwarranted interchange of embedding and
involution. For two vectors of this selected basis, every polynomial
coefficient belongs to $L^{\mathrm c}$; thus (FC.17a) proves that
their coefficient pairing, after embedding, is exactly (FC.15) with
the Hermitian order reversed as in (FC.16). Expand a general vector
$x=\sum_j a_jb_j$ only after those basis-pairing identities are proved.
Both sides are then the same complex quadratic form in the scalars
$s(a_j)$, with the usual complex conjugates. No equality
$s(a_j^*)=\overline{s(a_j)}$ is required for arbitrary coefficients
$a_j\in L$. This also explains the statement in any other coordinate
basis: transport the already obtained complex quadratic form by the
complex change-of-basis matrix.

The combined matrix of all basis values is injective. It has a left
inverse over $L$, whose embedded matrix remains a left inverse; hence
a nonzero complex coefficient vector has a nonzero class value.
The positive weighted Fock sum therefore makes the embedded Gram
matrix positive definite. The source's standard Hermitian frame
realizes precisely this Gram matrix. Polynomial injectivity (CV1)
identifies any proposed family $P_c$ in the norm statement with the
actual coefficient polynomials. Thus (FC.16)--(FC.17) prove the literal
\texttt{ipV\_norm\_tie}, rather than assuming a metric comparison.
The same argument in dimension $N_m-1$, with the terminal constant
weight shift, supplies \texttt{ipP} and \texttt{NormPTie} for the
terminal value construction of CV10. It proves no independent spectral
or metric estimate elsewhere in the reconstruction.

\paragraph{FC8: the period map is coefficient conjugation.}
Let $\mathcal C P$ be the frame function obtained by conjugating the
coefficients of $P$ and retaining the prescribed determinant twist.
It is conjugate linear and injective by CV1; it maps onto the prescribed
complex polynomial value space, because that space is defined over the
real integer highest-polynomial generators. For a rational unitary
matrix $\gamma$ and a character value $\chi$ of finite order,
$\gamma^{-1}=\gamma^*$ and $\overline{\chi^{-1}}=\chi$.
Substitution into the polynomial action gives the exact intertwining
identity
\[
 \mathcal C\bigl(\chi^{-1}\rho(\gamma)P\bigr)
          =\operatorname{tr}_{\gamma,\chi}(\mathcal C P).
 \tag{FC.18}
\]
Here $\mathcal C$ conjugates coefficients; it does not replace the
frame variable by its complex conjugate.

Evaluate an equivariant polynomial section on each preferred terminal
class representative and apply $\mathcal C$; call the resulting map
$j_M$. The source's \texttt{PerIn.H} is the whole group. Values on
preferred representatives therefore determine a section, proving
injectivity. Conversely, lift each prescribed class value uniquely by
$\mathcal C$ and extend by $\chi^{-1}\rho(\gamma)$ along its orbit.
Two possible moving elements differ by a stabilizer. Equation (FC.18),
the invariance of the class value and injectivity of $\mathcal C$
show that the lift is fixed by that stabilizer. The extension is
well-defined and equivariant. Thus $j_M$ maps onto exactly
the source's \texttt{ptC} space.

Conjugating the second polynomial in (FC.15) turns its Hermitian
coefficient pairing into the symmetric bilinear coefficient pairing.
Term by term in the weighted class sum this proves
\[
 \operatorname{period}(w,F)=\operatorname{ipP}(j_MF,w).
 \tag{FC.19}
\]
For nonzero $z\in\texttt{ptC}$, at least one class polynomial is
nonzero. Its Fock norm is positive, all class weights are positive,
and every other summand is nonnegative. Hence
$\operatorname{Re}\operatorname{ipP}(z,z)>0$. These statements prove
all five assertions of the literal \texttt{period\_conj} supplier
under the supplied column/period laws; they do not assert existence of
those laws for a new arithmetic host.

\paragraph{FC9: convert the chart levels and classes to Hermitian
carriers.}
This construction uses the auxiliary-prime descriptions
$\texttt{GLevH}$ of the same charts, constructed in FC14 below. At an auxiliary prime, a matrix
congruent to the identity lies in the chart level because its reduction
is the identity of its finite subgroup. Away from the auxiliary and
level primes the chart already has its stated full lattice stabilizer.
At a level prime, FC4 proves the remaining Hermitian condition exactly:
\[
 U_v=\{g\in\mathrm U(h):gL=L,\quad
                    (g-1)L^\#\subset L\}.            \tag{FC.20}
\]
In fact the old definition is the intersection with the stabilizer of
the nonzero prime-order functional, and FC4 shows that this is equivalent
to acting trivially on the whole discriminant module. This also explains
why one cannot simply replace a functional by an arbitrary zero map.

Use the coordinate basis of SF2 to carry the lattice, local groups and
global group back to the actual Hermitian space. Transport along the
single rank-index bijection takes addition, scalar multiplication,
duals, localizations and matrix conjugation to themselves. For each
class, use the very same $y_c$ and stabilizer homomorphism in these
coordinates. Their quotient bijection, lattice-transport condition,
injectivity and range formula are unchanged. This proves
\texttt{classTie\_of\_chart}, including every equality in
\texttt{SameCar}, rather than just an isomorphism of class sets.

At step zero the same transport produces \texttt{ColLevel}. Its
type-prime splitting and frames are the supplied $G$ and the same
\texttt{glev0.fr}. The additional claim that a rational stabilizer
preserves its class lattice follows as follows. Conjugating it by $y_c$
puts it in $U$, so it preserves every completed class lattice.
A full $\mathcal O$-lattice is the intersection of its finite localizations
inside $E^n$, by the denominator argument of LP1. It therefore preserves
the global lattice; apply the same argument to its inverse.

For $1\le i<m$, apply this construction to each chart fiber. At
$i=m$, the chart type is a singleton by FC3, so restriction to its
fiber and the inverse inclusion are exact inverse class maps. There
is no empty endpoint fiber. Use the residue equivalence (FC.3) and the
same place bijection for the remaining fields. These choices construct
the entire \texttt{CarH} output of \texttt{hermColsCh\_of\_forms},
conditional on the given standard forms and their level descriptions.
No separate Hermitian carrier existence assumption is needed after
those forms are available.

\paragraph{FC10: the step-zero value module is killed by the actual
uniformizer.}
The previous normalization proves this for the charts constructed here.
There is also a useful intrinsic proof for any supplied standard forms,
which is the source's \texttt{pi\_kills}. Choose a class, a vector in
its nonzero residue-functional image, and a lift $z$ using the
surjectivity of its projection. By \texttt{FvTie}, the corresponding
value $m=f(y_c^{-1}z)$ in the chart value group is nonzero.
Multiplication of $z$ by $\pi$ multiplies its projected value by
$\operatorname{res}(\pi)=0$. Linearity and \texttt{FvTie} therefore
give $\pi m=0$. An additive group of prime order is generated by
any nonzero element, so $\pi M=0$.

If a previously chosen \texttt{ColLevel} has the correct carriers but
different residue-module fields, transport the forms' functional along
the exact rank reindexing and use its value module. Nonzeroness,
$p$-torsion, order $p$, vanishing on the base lattice and its stabilizer
formula transport with it. This changes only those functional fields;
the lattice, local and global groups, class representatives, stabilizer
maps, splitting and frames remain identical. It proves the literal
\texttt{stepZeroHermitianLevel\_with\_formsFunctional} and
\texttt{stepZeroHermitianLevel\_uniformiser\_kills\_values} assertions.

\paragraph{FC11: what the naturality calculation does once a supported
move exists.}
This paragraph proves the algebraic part of the move comparison, with
its geometric premise exposed. Suppose an actual up/down/diamond move
keeps the chart label, has moving column-group element $g=(\gamma,h)$,
and has an adelic realization $y_c\delta k$ with $\delta$ equal to
the identity away from the specified level place. Let $c'$ be the
new preferred class. By \texttt{ReadsK},
\[
 \kappa=y_{c'}^{-1}\gamma^{-1}y_c\delta\in U,\qquad
 y_c\delta=\gamma y_{c'}\kappa,
 \qquad \chi_U(\kappa)=\chi_{\rm col}(g).              \tag{FC.21}
\]
The reason for membership is exact: $\gamma^{-1}$ moves the realized
point to the preferred point by the finite component $h$, which
preserves the chart coupling data. The level-reading equivalence of
FC3 applies, and its character-reading equality gives the last formula.
For frames with $z=\gamma z'$, the transported function is
$\chi_{\rm col}(g)f(\gamma^{-1}z)=\chi_U(\kappa)f(z')$.
This proves the character and frame parts of \texttt{upNatTie} and
\texttt{dnNatTieG\_cols}; the inverse $\kappa^{-1}$ in the source's
alternative factorization introduces exactly the displayed inverse
character, with no new scalar.

The elementary local-basis correction used in an up move is also exact.
If $a$ is a uniformizer at $P$ and $x-z\in PA$, then
$a^{-1}(x-z)\in A_P$, whence
\[
 (A+\mathcal Oa^{-1}x)_P=(A+\mathcal Oa^{-1}z)_P.
 \tag{FC.22}
\]
For a down-move functional, an identity known on the original global
lattice extends to its localization: choose $s\notin P$ making the
vector integral, use the residue-linearity of both sides, and cancel
$\bar s\ne0$. This is the cancellation in
\texttt{dn\_level\_functional}. These arguments do not construct
$\delta$, prove that the source's actual neighbor has the required
local lattice, or establish that a move preserves the coupling label.
Those are distinct geometric assertions retained below.

\paragraph{FC12: a common finite orbit carrier for the penultimate
charts.}
For a group $J$ and a finite family of finite-index subgroups
$U_\beta\le J$, set
\[
 Y=\coprod_\beta J/U_\beta,\qquad y_\beta=(\beta,U_\beta).
 \tag{FC.23}
\]
This is a finite $J$-set. Every point belongs to exactly one orbit
$J y_\beta$, and the stabilizer of $y_\beta$ is exactly $U_\beta$.
If $\chi_\beta$ is the restriction of a common character $\chi_J$,
its stabilizer character is also the required one. These assertions
follow directly by multiplying cosets and retaining their disjoint-union
label. They prove the actual \texttt{exists\_orbits} part of
\texttt{stepBefore\_oneLevel}.

For the actual penultimate charts, the supplied index law
\texttt{last\_uniq} makes the finite column-group index independent
of the class. Choose a point whose coupling coset contains zero.
Such a point exists by taking the zero coset of the prescribed coupling
lattice. Membership of zero is preserved by every linear column action,
so its preferred representative still has zero coupling reduction;
let $\beta_0$ be its chart. All zero-containing coupling cosets have
the same reduction family: translating their defining lattice by zero
leaves that lattice, and a reduction frame sends it to the prescribed
standard lattice. In the reference case these reductions at the auxiliary
primes are the singleton zero.

Here is the exact group calculation. Let $J_0$ be the unitary integral
adelic group whose reduction fixes the last row at every selected level
prime, and let $r$ be its product of auxiliary-prime reduction maps.
The AD1 base point normalizes $J_0$: if
$e_n^{\mathsf T}k^{-1}=b e_n^{\mathsf T}$ after reduction, then
$b$ is a unit (multiply by $k$), and the equation
$e_n^{\mathsf T}x=e_n^{\mathsf T}$ is equivalent to the same
equation for $k^{-1}xk$. At auxiliary primes $k=1$, so conjugation
does not change $r(x)$. The point-reading argument of FC3 consequently
gives
\[
 U_\beta=\{x\in J_0:r(x)^{-1}\in H,
                   \ r(x)^{-1}C_\beta=C_\beta\},
 \qquad \chi_\beta(x)=\theta(r(x)^{-1}),             \tag{FC.23a}
\]
where $H$ and $\theta$ are the same finite column group and character
for all $\beta$. The zero reduction $C_{\beta_0}$ is fixed by $H$.
Thus choosing $J$ to be the already constructed chart level at
$\beta_0$ gives $U_\beta\subset J.U$, equality of all base
lattices, and restriction of the character of $J$ to $U_\beta$.
Moreover $\ker r\subset U_\beta$ and $r$ has finite image, so
each $U_\beta$ has finite index in $J.U$. Applying (FC.23) now
constructs the actual finite orbit carrier with its exact stabilizers.

This proves the common level, inclusions, character restrictions and
orbit part of \texttt{stepBefore\_oneLevel}. What remains there is
the specified identification of its reduction group with
$\texttt{HSym}$ and its character with the displayed
$\texttt{lamSym}$ formula. The generic reduction-image construction
of FC14 alone does not prove that particular matrix shape and character
formula.

\paragraph{FC13: construct the one common auxiliary-prime package.}
Write $q_j=g_j$ for the two distinct supplied auxiliary rational primes,
and $a_j=\texttt{piG}(j)$, with $a_j\bar a_j=q_j$.
The supplied complete splitting at these primes gives, for each
$v\mid q_j$, a splitting
$E\otimes_FF_v\simeq F_v\times F_v$ which carries conjugation
to exchange of factors and scalars to the diagonal. Its residue field
is the prime field $\mathbb F_{q_j}$. Orient the splitting so that the
first component of $\bar a_j$ has valuation less than one.
This is always possible: the two components of $a_j$ have product
$q_j$, whose valuation is less than one. If the required component is
not small, the other is; exchange the two factors. Both displayed
splitting identities survive that exchange. The equality
$a_j\bar a_j=q_j$ and the chosen $a_j$ themselves are unchanged.

Choose one selected auxiliary prime $w_j$ of $E$ above $q_j$; its
existence follows from the supplied positive cardinality of the selected
set. The rings $k_v$ and $\mathcal O/w_j$ both have order $q_j$.
The unique unital map from $\mathbb F_{q_j}$ to either is an
isomorphism: its kernel is zero and domain and codomain have the same
finite cardinality. Let $e_v:\mathbb F_{q_j}\simeq k_v$ and
$e_w:\mathbb F_{q_j}\simeq\mathcal O/w$ denote these maps.
For the coefficient-field embedding $\sigma_0:L\hookrightarrow\mathbb C$,
define
\[
 \psi_{j,v}(x)=\sigma_0\!\left(
       \psi^{\rm col}_{w_j}(e_{w_j}e_v^{-1}x)\right).
 \tag{FC.24}
\]
It is an additive character. If it were trivial, surjectivity of the
field isomorphisms and injectivity of $\sigma_0$ would make the supplied
column character trivial. Thus it is nontrivial. The column-character
law says that $\psi^{\rm col}_w$ has the same value on every integer
at every selected $w\mid q_j$. Since every element of these prime
fields is an integer residue, (FC.24) agrees with the required character
at every such $w$, not only at $w_j$.

The supplied multiplicative-character law gives an injective homomorphism
$\omega_j:\mathbb F_{q_j}^{\times}\to L^{\times}$ for which
the $s$th column character is $\omega_j^{2(s+1)}$.
Extend $\sigma_0\omega_j$ to a multiplicative character on the field,
and transport its powers by $e_v^{-1}$:
\[
 \chi_{j,v,b}(x)
       =\bigl(\sigma_0\omega_j\bigr)^{2(b+1)}(e_v^{-1}x).
 \tag{FC.25}
\]
This character has the prescribed values on every nonzero integer
residue. Injectivity gives the primitive character
$\sigma_0\omega_j$ order $q_j-1$, since
$\mathbb F_{q_j}^{\times}$ is cyclic of that order. If the powers
in (FC.25) for $a,b<m$ were equal, transporting them to $w_j$ would
make the corresponding supplied column characters equal. Their
injectivity law gives $a=b$. Equations (FC.24)--(FC.25) also show
agreement across all $v\mid q_j$ on integer residues and exhibit the
single primitive multiplicative character required by
\texttt{CharacterCompatible}.

The oriented splittings, the original $a_j$, and these characters are
the complete $G=\texttt{GLoc}$ object. They prove all three fields of
\texttt{IdG} and the complete \texttt{CharacterCompatible} assertion
in \texttt{exists\_gLoc\_idG\_chiTied}. Choose this $G$ once,
independently of the word $u$. Since $G.\texttt{piG}=X.\texttt{piG}$,
the required membership of its uniformizers in the designated subfield
is exactly the supplied $X.\texttt{piG\_K}$ clause. Complete splitting,
the original character laws and that subfield clause remain part of
the host data, not consequences of this transport argument.

\paragraph{FC14: construct the chart reduction groups and their
characters.}
Fix any chart $S$ from FC3 and any $v\mid q_j$.
Because $q_j\ne p$, the element $p$ is a unit at $v$. Its base
lattice satisfies
$p\mathcal O^n\subset L_\pi\subset\mathcal O^n$, so the
completed lattice is the standard integral lattice. This is the explicit
denominator calculation of SF5. The first split projection therefore has
the identity matrix as a frame, including the required existential
lifting of every integral first-projection vector.

Let $V_v$ be the unitary stabilizer of this completed lattice and let
$r_v:V_v\to\mathrm{GL}_n(k_v)$ be reduction of the first split
matrix. Its second split matrix is the inverse transpose of the first,
because the Hermitian Gram matrix is $I$. Hence $r_v(x)=I$ implies
that both split matrices of $x$ and $x^{-1}$ are congruent to $I$.
Embed $x$ at the single place $v$, with identity elsewhere. It preserves
the completed lattice and every reduction frame of the decorated point:
at $v$ this follows from the congruence, and elsewhere it is the identity.
The base point $k_u$ is also the identity at $v$.
Thus right multiplication by this embedded element leaves the same
decorated point and its coupling reductions. The exact \texttt{ReadsK}
level and character clauses give
\[
 \ker r_v\subset U_v,\qquad
 \chi_U(\iota_v x)=1\quad(x\in\ker r_v),             \tag{FC.26}
\]
where $\iota_v$ denotes the supported embedding.

Define $H_v=r_v(U_v)$ and
$\lambda_v(r_vx)=\chi_U(\iota_vx)^{-1}$ for $x\in U_v$.
Equation (FC.26) makes this independent of the lift and multiplicative.
Moreover
\[
 U_v=r_v^{-1}(H_v)\quad\hbox{inside }V_v.             \tag{FC.27}
\]
Indeed, if $r_vx=r_vy$ with $y\in U_v$, then
$y^{-1}x\in\ker r_v\subset U_v$.
Only finitely many places divide the two distinct nonzero $q_j$.
For $x\in U$, remove its components at these places using the
$\iota_v$. The remaining element is the identity at every bad place;
the given \texttt{StepLevel} character is trivial on that subgroup.
Consequently
\[
 \chi_U(x)=\prod_{j=0}^1\prod_{v\mid q_j}
                  \lambda_v(r_vx_v)^{-1}.            \tag{FC.28}
\]
Each $H_v$ is finite, so each $\lambda_v$ has finite image and all
its values have absolute value one. Equations (FC.27)--(FC.28), with
the identity frame, supply every field of \texttt{GLevH}.
This proves the actual \texttt{chart\_gLevH} supplier from the
constructed charts. No further auxiliary-prime description is assumed
for those charts.

\paragraph{FC15: the explicit step-zero level and character.}
At $i=0$, $n=N_0=k$, the coupling set is empty, and the supplied
index law identifies the original finite column subgroup at each
selected auxiliary prime with the upper unitriangular group
$\mathrm U_n(k_v)$. This identification is part of
$\texttt{idx\_ok.zero\_G}$; its matrix content is simply
$A_{ab}=0$ for $b<a$ and $A_{aa}=1$. The source's
$\texttt{InJ}(k,t,N_0,\mathrm{false})$ has exactly this content:
all its block labels are singleton indices, while its extra fixed-column
and last-row conditions have no applicable index.

The oriented first split factor in FC13 corresponds to the unique
selected prime $w\mid\bar a_j$ above $v$. Its residue map
$\rho_w:\mathcal O/w\simeq k_v$ is the map induced by that
first projection. It carries every integer to itself. Reading the
same point after a matrix supported at $v$ changes its reduction frame
by the inverse matrix at $w$ and leaves all other selected reductions
unchanged. Since $\mathrm U_n$ is a group, the inverse has the
required shape exactly when the matrix itself does. The absence of
coupling and twist-prime data at step zero, together with \texttt{ReadsK},
therefore proves
\[
 x\in U_v\quad\Longleftrightarrow\quad
 r_v(x)\in\mathrm U_n(k_v)\qquad(x\in V_v).          \tag{FC.29}
\]
This proves the specified group identity, rather than merely choosing
the image subgroup as in FC14.

For $A\in\mathrm U_n$, put
$s(A)=\sum_{a=1}^{n-1}A_{a,a+1}$. Matrix multiplication gives
$s(AB)=s(A)+s(B)$: in the $(a,a+1)$ entry only the two diagonal
terms survive. Thus $s(A^{-1})=-s(A)$.
At step zero the column character is
$\prod_w\psi_w^{\rm col}(s(A_w))$; there are no determinant
block factors or twist-prime factors. The chart reads
$A_w=r_w(x)^{-1}$. Applying (FC.24), the actual residue maps
$\rho_w$, and the last identity gives
\[
 \chi_U(x)=\prod_{j=0}^1\prod_{v\mid q_j}
               \psi_{j,v}\bigl(s(r_vx_v)\bigr)^{-1}. \tag{FC.30}
\]
This is exactly the source's \texttt{lamJ} formula with
$(a,nb,skip,fc,zl)=(k,0,k,N_0,\mathrm{false})$.
The finite-product reindexing is legitimate: restriction is a bijection
on each selected auxiliary-prime set, by the same conjugate-prime and
coprimality argument as FC1, and the two sets are disjoint because
$q_0\ne q_1$. The additive-character values have finite order and
absolute value one. The identity frame, (FC.29), (FC.30), and this norm
calculation prove all fields of \texttt{glev\_zero}, including
\texttt{fr0}$=I$.

\paragraph{FC16: the exact joins to the comparison hypotheses.}
The level-zero coefficient-field condition has a short proof which is
stronger than requiring the particular construction (FC.24).
For any additive character $\psi:k_v\to\mathbb C$ at $v\mid q_j$,
\[
 \psi(x)^{q_j}=\psi(q_jx)=\psi(0)=1.                 \tag{FC.31}
\]
The supplied \texttt{Base.roots\_Ks} contains a primitive
$q_j(q_j-1)$st root $\zeta$ in $L$, with its specified complex
embedding. Every complex $q_j$th root of unity is a power of
$\sigma_0(\zeta)$: it is a root of
$X^{q_j(q_j-1)}-1$, whose distinct roots are precisely those powers.
Hence $\psi(x)=\sigma_0(\zeta^a)$ for some integer $a$.
Applying this to any $G.\texttt{psi0}$ proves the literal
\texttt{LevelZeroCharacterRational} condition. This proves the actual
\texttt{additiveCharacter\_typePrimes\_values\_coefficientField}
supplier without using its \texttt{IdG} argument, and without another
global forms assumption.

The step-zero carrier construction in FC9 uses only the step-zero
chart, its lattice/stabilizer data and the level of FC15. Call this
level $L_0(u)$. Any completion of our partial forms data to
$T_u:\texttt{Forms0}$, keeping these specified fields, satisfies
\texttt{Joint0} with $L_0$ directly: its five assertions are equality
of the reindexed base lattice, local groups, global group, class matrices
and stabilizer matrices, exactly the \texttt{car0} identities in FC9.
They are unaffected by FC10's functional replacement.

This also removes a need to supply separate companion ideals for
these carriers. The FC3 base-lattice equality and FC9's exact reindexing
give $L_0(u)=L_{\pi_u}$ in the standard space. SF2 supplies
\texttt{StdTied}, and SF3 gives
\[
 L_0(u)+(\bar\pi_u)L_0(u)^\#=\mathcal O^{N_0}.
 \tag{FC.32}
\]
Thus the comparison ideals are the explicit $(\bar\pi_u)$ for these
same level-zero carriers. Their construction needs no terminal
\texttt{CarT} data. Once the remaining fields of \texttt{Forms0}
are supplied, VE9's \texttt{StdTied}, \texttt{Joint0},
\texttt{LevelZeroCharacterRational} and companion premises are all
discharged by these exact joins. FC19 supplies the full
\texttt{Forms0} directly, using NB7--NB8 for its naturality fields.

\paragraph{FC17: the distinguished-block lattice comparison.}
Fix a coefficient prime $w$ above $p$ which is seen, and a selected
level prime $P$. For a class decoration, put $A=L^\#$ and let
$K_P\subset A$ be the kernel of its $P$-residue functional.
The decorated-lattice laws and the denominator argument of FC5 give
\[
 (K_P)_P=(L+PA)_P,\qquad (K_P)_Q=A_Q
             \quad(Q\ne P\text{ selected}).          \tag{FC.33}
\]
For the second equality, choose $a\in P\setminus Q$.
Multiplication by $a$ kills the $P$-residue of every vector in $A$,
while $a$ becomes a unit at $Q$; this proves the nontrivial inclusion.
For the first, zero residue is exactly local membership in $L$,
and the kernel after clearing denominators is $L+PA$, as in (FC.12).

Choose rational $p$-local basis matrices $B,B'$ for $A$ and
$L+PA$ respectively. Such bases exist because both are full lattices;
for the second lattice use $L\subset L+PA\subset A$.
The rational $p$-local basis condition implies the prime-local basis
condition at every $Q\mid p$, by CV6. Indeed a rational integer
prime to $p$ cannot belong to $Q$, and a product of the finitely many
coordinate denominators clears the converse implication.
At a seen $w$ the supplied place correspondence is exact:
the place below $w$ through the block embedding labelled
$\tau_Q=\texttt{tauAt}(w,Q)$ is $Q$.
Thus (FC.33) gives a basis of $K_P$ in every block, namely
\[
 B_\tau=\begin{cases}B',&\tau=\tau_P,\\B,&\tau\ne\tau_P.
            \end{cases}                              \tag{FC.34}
\]
Apply the exact localized Weyl-lattice transport of CV7 to this family.
Membership in the $w$-localization of the polynomial lattice attached
to $K_P$ is equivalent to one integral-span polynomial evaluated with
$B_\tau^{-1}$ in each block. Formula (FC.34) is precisely the
piecewise inverse-frame formula in \texttt{BwTieG}. CV8 takes the
product of the classwise denominators to turn these conditions into
membership in the global $\texttt{Bw}$ localization. All frame maps,
embeddings and the determinant shift are the same identities as CV9.
The source quantifies over every pair $B,B'$; the forward implication
works for every such pair, and the reverse uses the existence just proved,
so it is not vacuous. This proves the full actual \texttt{bw\_tie}
supplier. These local matrices describe the class lattice and do not
replace the common global comparison frame of SF1--SF6. The assertion
is restricted to seen $w$ of residue characteristic $p$, as in the
source; it makes no claim about unseen coefficient primes.

\paragraph{FC18: the actual diamond move from its defining functional
equation.}
Let $s\in k_P^\times$, choose an integer $a$ representing $s$,
and let $\delta$ satisfy the source's \texttt{IsDia} condition at
$v=P\cap\mathcal O_F$. This means that $\delta$ is the identity
away from $v$, is unitary, preserves the base lattice, and satisfies
\[
             a f(\delta x)=f(x)\qquad(x\in L^\#).
 \tag{FC.35}
\]
Its inverse preserves the same lattices. FC4 therefore proves that
both matrices have integral entries. Substitute $\delta^{-1}x$
in (FC.35), and use FC6's exact functional identification, to obtain
\[
 \rho_P\bigl(\bar\pi(\delta^{-1}x)_n\bigr)
                =s\,\rho_P(\bar\pi x_n).
 \tag{FC.36}
\]
Every coordinate vector belongs to $L^\#$, and $\bar\pi$ is a
unit at $P$. Testing (FC.36) on them gives
\[
        (\delta^{-1})_{nj}\equiv s\,\delta_{nj}\pmod P.
 \tag{FC.37}
\]
Here the symbol on the right is the Kronecker delta.
Thus the source's integral-support-last-row conclusion follows directly
from the normalized functional; no choice of a special diamond matrix
is needed.

Suppose $y_ck_u$ realizes the preferred decorated point. AD1 gives
$k_u^{-1}$ last row $\bar\pi e_n^{\mathsf T}$ modulo $P$,
and $k_u$ preserves the standard integral adelic lattice.
For a vector in the original decorated lattice, $y_c^{-1}z$ is
therefore integral. Reducing the product
$k_u^{-1}\delta^{-1}y_c^{-1}z$ is legitimate. Equations (FC.37)
and the last-row formula multiply its old functional by $s$.
The underlying lattice is unchanged, because $\delta$ and its inverse
are integral. Every auxiliary reduction, every other selected level-prime
functional, and every twist-prime frame is unchanged, since its place
is outside $v$. No selected prime is conjugate to $P$.
Consequently $y_c\delta k_u$ realizes exactly the source's
\texttt{dia1} decoration.

This operation changes only the $P$-functional. In particular its
coupling reductions are unchanged. The supplied level law
\texttt{dia\_act} says that the operation commutes with the column
group action, so the sets of attainable coupling reductions are also
unchanged. The preferred-normalization argument of FC3 proves that the
new preferred class has the same chart label. Apply FC11 to the
realizer just constructed: it gives the exact adelic factorization,
character scalar and frame transport. This proves all clauses of the
actual \texttt{dilTied} supplier, under the literal supplied column
laws and the chart realizability input. It also supplies the dilation
field required when CV9's chart values are made into an endpoint
\texttt{ValTied} record.

\paragraph{FC19: direct construction of Forms0 and its matching comparison data.}
Fix the actual $D=I.\texttt{rec7}\,l\,Z$ and the supplied
\texttt{Ok}, \texttt{Pins}, \texttt{Laws} and \texttt{Later} data.
Retain the explicit initial realization hypothesis used in FC3 and
NB6: every valid framed decorated point needed there has a unitary
finite-adelic realizer in the full \texttt{IsPtOf} sense.
We now construct, for every word $u$, an actual
$T_u:\texttt{Forms0}\,D\,u$, level data $L_0$, and all of
VE9's global joins. The construction uses only step-zero form data.

\emph{First construct the form record.}
Choose the common $G$ of FC13 once for all words. At $i=0$, FC3
gives a unique chart label $\beta_u$ and
$\operatorname{chOf}(c)=\beta_u$ for every class $c$. FC4 normalizes
its functional without changing its level, class matrices or stabilizer
matrices. Define $\operatorname{cls0}_u$ by retaining this level and
setting
\[
 y_u(c)=K_u(\beta_u).y(c,\operatorname{chOf}(c)=\beta_u),\qquad
 s_u(c,\gamma)=K_u(\beta_u).\operatorname{stabI}
                      (c,\operatorname{chOf}(c)=\beta_u)(\gamma).
 \tag{FC.38}
\]
The map from the original class set to this chart fiber is a bijection
with inverse the first projection. Thus the class bijection, carried
lattice identity, and injective stabilizer with its exact image are
the same assertions as FC3. This is \texttt{clsOfChart} itself,
with no replacement representatives.

Take $\operatorname{atT}_u$ from FC1 and the frame equivalence $e_Y$
from SF2. The determinant witness in \texttt{ValTied.de} is $C=I$:
the Gram matrix is $I$, the integral-coordinate lattice is
$\mathcal O^{N_0}$, and the original determinant label equals
$\det(e_Y z)$ at every embedding. CV9 supplies the complex-form,
coefficient-field and localized integral-value identities.
Because the chart set is a singleton, a function satisfies its
chartwise adelic condition if and only if its restriction through
the bijection in (FC.38) does. Surjectivity of that bijection proves
both implications by evaluation on every class; no extension is lost.

The global integral field of \texttt{ValTied} also follows from
these localized assertions. For any $\mathcal O_L$-submodule $M$
of an $L$-vector space and any $x$ in every localization $M_w$,
put $\mathfrak b_x=\{a\in\mathcal O_L:ax\in M\}$.
If $x\notin M$, this ideal is proper and lies in a maximal ideal
$w$. Membership $x\in M_w$ supplies an $a\notin w$ with
$ax\in M$, contradicting $\mathfrak b_x\subseteq w$.
Every maximal ideal of $\mathcal O_L$ is a coefficient prime, so
\[
                         M=\bigcap_w M_w.           \tag{FC.39}
\]
Here localization membership supplies such a denominator by taking
the product of the finitely many denominators in an expression for
$x$ in the localized span. This argument does not require an
unproved freeness statement for $M$.
Apply it to the actual form lattice, using the full prime correspondence
and CV9 at every coefficient prime. FC7 supplies the exact inner-product
field and FC18 supplies both dilation fields. These data are precisely
\texttt{valOfG}, and give $\operatorname{val0}_u$ with unchanged $e_Y$.

The remaining fields of \texttt{Forms0} are now specified individually.
Its \texttt{gl} is the same $G$, and \texttt{gl\_K} follows from
$G.\texttt{piG}=X.\texttt{piG}$ and the supplied
$X.\texttt{piG\_K}$. Its \texttt{glev0} is the actual level of FC15,
with parameters
$(a,t,nb,skip,fc,zl)=(k,t,0,k,N_0,\mathrm{false})$ and frame $I$.
FC5 gives \texttt{lv0} by retaining its quotient projections;
FC6 gives the additive residue isomorphisms \texttt{fMe0} and
their exact \texttt{fv0} identity. NB7 and NB8 give
\texttt{up\_natural\_level\_zero} and
\texttt{down\_natural\_level\_zero}: substitute the unique chart
label, discard the now automatic chart-equality clause, and retain
the same $\gamma,\eta,\kappa$, lattices, functionals and frame scalar.
These are the source's \texttt{upOfG} and \texttt{dnOfG} conversions.
FC7 gives \texttt{norm0}, and FC17 gives \texttt{bw0}; their endpoint
conversions merely use the identical $e_Y$. We have therefore filled
every field of \texttt{Forms0} directly. Neither the period map nor
the penultimate or terminal form records are needed for this step.

\emph{Next construct the matching level carrier.}
Let $r:k=N_0$ be the fixed \texttt{rank\_eq\_first\_level}
identity. Apply the reindexing construction in FC9 to
$\operatorname{cls0}_u.\operatorname{lev}$, its actual class and
stabilizer matrices, $G$, and the just constructed
$\operatorname{glev0}_u$. This construction is the proof of
\texttt{exists\_colLevel} with only its step-zero arguments:
every occurrence of its apparent full form parameter is through
\texttt{cls0}, \texttt{gl}, or \texttt{glev0}.
Transport the lattice, local groups, global group, functional,
class matrices and stabilizer matrices by $r$; transport every
relation by the same coordinate permutation. The upper-unitriangular
description in FC15 gives the required auxiliary level clause for
\texttt{ColLevel}; the finite prime of $F$ above $q_0$ required in
that record exists by lying over. The stabilizer preserves each actual global
class lattice because it preserves every completed lattice. For a rational
vector, membership in the completion of a prime-local lattice is the same
as membership in that localization: write it in a local basis and test
that each rational coefficient has nonnegative valuation. LP1's local
membership criterion then recovers the global lattice.
Thus these data satisfy all \texttt{ColLevel} fields. Call the result
$L_0(u)$ and choose it for each word.

In particular its five carried objects satisfy exactly
\[
\begin{split}
 \operatorname{reL}_{r}(L_0(u).L_b)&=
                         \operatorname{cls0}_u.\operatorname{lev}.L_b,\\
 \operatorname{reGL}_{r}(L_0(u).U_{\rm loc,v})&=
                         \operatorname{cls0}_u.\operatorname{lev}.U_{\rm loc,v},\\
 \operatorname{reGL}_{r}(L_0(u).U)&=
                         \operatorname{cls0}_u.\operatorname{lev}.U,\\
 \operatorname{reGL}_{r}(L_0(u).y_c)&=y_u(c),\qquad
 \operatorname{reGL}_{r}(L_0(u).s_{c,\gamma})=s_u(c,\gamma).
 \tag{FC.40}
\end{split}
\]
For groups the notation denotes their image under the reindexing
homomorphism. Equation (FC.40) is the literal
$\texttt{Joint0}\,L_0\,u\,T_u$, not merely an abstract isomorphism
of class spaces. The splitting used in $L_0(u)$ is $G.\texttt{sp}$,
and its local frame reindexes to the same identity frame of
$\operatorname{glev0}_u$. Functional normalization has not changed
any term of (FC.40).

Finally SF2 proves \texttt{StdTied} for this actual $D$, including
all of its positive spaces and its terminal positive space; it is
independent of the missing terminal incidence construction. FC16 gives
\texttt{LevelZeroCharacterRational} for $T_u.\texttt{gl}=G$.
The transported base lattice is exactly
$L_{\pi_u}=\mathcal O^{k-1}\oplus\pi_u\mathcal O$ because the
single rank reindexing retains the last coordinate. Put
$\mathfrak a_u=(\bar\pi_u)$. SF3 then proves
\[
 L_0(u).L_b+\mathfrak a_u\bigl(L_0(u).L_b\bigr)^\#
       =\mathcal O^k
       =D.\operatorname{sp}.\operatorname{crdLat0}
                         \bigl((D.\operatorname{sp}.X_0).L_0\bigr).
 \tag{FC.41}
\]
The common standard comparison frame can be taken to be $B=I$.
Consequently the \texttt{Forms0}, \texttt{Joint0}, \texttt{StdTied},
level-zero character-rationality and companion hypotheses in VE9 all
have actual simultaneous suppliers here, under the stated true-column
and initial realizability hypotheses. No full \texttt{CarT},
\texttt{FormsTiedSym}, or \texttt{ColsTiedSym} has been assumed in
order to construct these comparison inputs.

\emph{The separate full-column boundary.}
Ordinary choice over the word/step sets also assembles the other
pointwise fields proved in FC1--FC18. Use the same $G$ everywhere;
unused steps $i>m$ may have singleton chart type. CV9 and FC7 give
the chart values and norm ties, FC8 gives the period tie, and CV10
and FC7 give terminal row values and norms once its row chart is
supplied. The full simultaneous \texttt{FormsTiedSym} route still
requires the particular penultimate $\texttt{HSym}/\texttt{lamSym}$
identification (FC12 already gives the actual common level and finite
carrier), the terminal row level and \texttt{ReadsP}, and the terminal
vector/incidence and adelic evaluation assertion \texttt{til\_ties}.
The corresponding full \texttt{CarT} record must be tied to the same
objects before \texttt{colsTiedSym\_std} is closed. Initial decorated
realizability and the existence of the original arithmetic host,
characters, column laws and later-stage data remain explicit inputs.
